\documentclass[../DoAn.tex]{subfiles}
\begin{document}

Chương này trình bày kết quả khảo sát hiện trạng các giải pháp theo dõi vị trí trẻ em hiện có trên thị trường, phân tích tổng quan chức năng hệ thống thông qua các biểu đồ use case, đặc tả chi tiết các use case quan trọng nhất, và xác định các yêu cầu phi chức năng cần đáp ứng.

\section{Khảo sát hiện trạng}
\label{section:2.1}

Trong những năm gần đây, cùng với sự phát triển của công nghệ Internet of Things (IoT), hệ thống định vị vệ tinh toàn cầu (GNSS) và mạng truyền thông di động, nhiều giải pháp theo dõi vị trí trẻ em đã được nghiên cứu và thương mại hóa. Dựa trên phương thức hoạt động và công nghệ sử dụng, các giải pháp hiện nay có thể được chia thành ba nhóm chính: ứng dụng theo dõi trên điện thoại thông minh, thiết bị định vị ngắn tầm dựa trên Bluetooth/UWB và các thiết bị GPS Tracker chuyên dụng sử dụng mạng di động.

\textbf{Nhóm 1 – Ứng dụng theo dõi trên điện thoại thông minh.} Life360 là một trong những ứng dụng tiêu biểu của nhóm này, cho phép các thành viên trong gia đình chia sẻ vị trí với nhau thông qua GPS của điện thoại thông minh \cite{life360}. Ứng dụng cung cấp nhiều tính năng như theo dõi vị trí theo thời gian thực, thiết lập vùng an toàn địa lý (Geofence) và gửi thông báo khi người dùng đi vào hoặc ra khỏi vùng đã thiết lập. Tuy nhiên, giải pháp này phụ thuộc vào việc trẻ phải luôn mang theo điện thoại thông minh, thiết bị phải được kết nối Internet và ứng dụng phải hoạt động liên tục. Điều này không phải lúc nào cũng khả thi đối với trẻ nhỏ hoặc trong những trường hợp nhà trường hạn chế học sinh sử dụng điện thoại. Bên cạnh đó, điện thoại thông minh hoặc đồng hồ thông minh còn tích hợp nhiều chức năng như trò chơi, mạng xã hội và các ứng dụng giải trí, dễ khiến trẻ bị phân tâm, làm giảm sự tập trung trong học tập và sinh hoạt hằng ngày.

\textbf{Nhóm 2 – Thiết bị định vị ngắn tầm.} Apple AirTag và Google Find My Network là những đại diện phổ biến của nhóm này. Các thiết bị sử dụng công nghệ Bluetooth Low Energy (BLE) kết hợp Ultra-Wideband (UWB) để hỗ trợ định vị với độ chính xác cao trong phạm vi gần. Ưu điểm của giải pháp là thiết bị có kích thước nhỏ gọn, tiêu thụ năng lượng thấp và không yêu cầu kết nối mạng di động. Tuy nhiên, việc xác định vị trí phụ thuộc vào mạng lưới các thiết bị Apple hoặc Android ở gần để chuyển tiếp tín hiệu, do đó độ tin cậy giảm đáng kể tại các khu vực thưa dân cư hoặc ít thiết bị tham gia mạng lưới. Ngoài ra, cơ chế này không hỗ trợ cập nhật vị trí liên tục theo thời gian thực như các thiết bị sử dụng GPS kết hợp mạng di động.

\textbf{Nhóm 3 – GPS Tracker chuyên dụng.} Các thiết bị như Optimus 2.0 hay Bouncie tích hợp bộ thu GPS và mô-đun truyền thông di động để truyền dữ liệu vị trí về máy chủ theo thời gian thực. Nhóm giải pháp này đáp ứng tốt yêu cầu theo dõi vị trí liên tục và hoạt động độc lập mà không cần trẻ mang theo điện thoại thông minh. Tuy nhiên, phần lớn các sản phẩm thương mại yêu cầu người dùng trả phí dịch vụ hằng tháng, phụ thuộc vào nền tảng của nhà sản xuất, khả năng tùy chỉnh các chức năng như Geofence còn hạn chế và thường không cung cấp API mở để tích hợp hoặc mở rộng hệ thống theo nhu cầu riêng.

Bảng \ref{table:comparison} tổng hợp và so sánh các nhóm giải pháp theo các tiêu chí quan trọng.

\begin{table}[H]
\centering
\caption{So sánh các giải pháp theo dõi vị trí trẻ em hiện có}
\label{table:comparison}
\begin{tabular}{|l|c|c|c|c|}
\hline
\textbf{Tiêu chí} & \textbf{Life360} & \textbf{AirTag} & \textbf{GPS Tracker TM} & \textbf{Đề tài này} \\
\hline
Không cần điện thoại trẻ    & Không & Có       & Có        & Có      \\
\hline
Theo dõi thời gian thực     & Có    & Không    & Có        & Có      \\
\hline
Geofence tự động            & Có    & Không    & Hạn chế   & Có      \\
\hline
Cảnh báo đa kênh            & Có    & Không    & Hạn chế   & Có      \\
\hline
Chi phí vận hành            & Miễn phí/Thuê bao & Một lần & Thuê bao & Thấp (SIM) \\
\hline
Khả năng tùy chỉnh          & Thấp  & Rất thấp & Thấp      & Cao     \\
\hline
\end{tabular}
\end{table}

Từ phân tích trên, hệ thống cần xây dựng phải giải quyết được các điểm yếu mà các giải pháp hiện tại chưa đáp ứng đầy đủ: hoạt động độc lập không cần điện thoại của trẻ, theo dõi liên tục theo thời gian thực, geofencing linh hoạt, cảnh báo đa kênh, và chi phí vận hành thấp.

\section{Tổng quan chức năng}
\label{section:2.2}

\subsection{Biểu đồ use case tổng quát}
\label{subsection:2.2.1}

Hệ thống có ba tác nhân chính. \textbf{Phụ huynh} là người dùng chính của ứng dụng di động, thực hiện các thao tác quản lý và theo dõi. \textbf{Thiết bị IoT} (thiết bị GPS đeo trên trẻ) là tác nhân tự động, định kỳ đọc tọa độ GPS và gửi lên server mà không cần tương tác người dùng. \textbf{Dịch vụ thông báo} (Telegram Bot) là tác nhân bên ngoài nhận lệnh gửi cảnh báo từ backend khi phát hiện vi phạm geofence hoặc tín hiệu SOS.

Các use case chính của hệ thống gồm: (1) Đăng ký tài khoản; (2) Đăng nhập hệ thống; (3) Quản lý thiết bị GPS; (4) Xem vị trí trẻ theo thời gian thực trên bản đồ; (5) Quản lý vùng an toàn địa lý (geofence); (6) Cảnh báo vi phạm vùng an toàn; (7) Gửi tín hiệu SOS khẩn cấp; và (8) Xem thống kê hành trình (quãng đường, số lần vi phạm, lịch sử trên bản đồ).

\begin{figure}[H]
    \centering
    \includegraphics[width=0.85\textwidth]{Hinhve/usecase_tong_quat.png}
    \caption{Biểu đồ use case tổng quát của hệ thống}
    \label{fig:usecase_total}
\end{figure}

\subsection{Biểu đồ use case phân rã – Quản lý thiết bị GPS}
\label{subsection:2.2.2}

Use case \textit{Quản lý thiết bị GPS} được phân rã thành bốn use case con tương ứng với bốn thao tác cơ bản: đăng ký thiết bị mới bằng cách nhập mã thiết bị để liên kết với tài khoản, xem danh sách các thiết bị đã đăng ký cùng trạng thái kết nối, đổi tên thiết bị để dễ nhận biết, và xóa thiết bị khỏi tài khoản khi không còn sử dụng.

\begin{figure}[H]
    \centering
    \includegraphics[width=0.85\textwidth]{Hinhve/usecase_phan_ra.png}
    \caption{Biểu đồ use case phân rã – Quản lý thiết bị GPS}
    \label{fig:usecase_phanra}
\end{figure}

\subsection{Quy trình nghiệp vụ}
\label{subsection:2.2.3}

Hệ thống có ba quy trình nghiệp vụ cốt lõi được mô tả qua các biểu đồ hoạt động dưới đây.

Quy trình cập nhật vị trí theo thời gian thực (Hình \ref{fig:ad01}) thể hiện luồng dữ liệu định kỳ từ thiết bị IoT đến ứng dụng di động: thiết bị đọc tọa độ GPS mỗi 5 giây và gửi lên MQTT Broker, backend nhận bản tin, lưu vào cơ sở dữ liệu và đẩy vị trí mới đến ứng dụng qua WebSocket, ứng dụng cập nhật marker trên bản đồ.

\begin{figure}[H]
    \centering
    \includegraphics[width=0.85\textwidth]{Hinhve/activity_vitri.png}
    \caption{AD01 – Biểu đồ hoạt động: Cập nhật vị trí GPS thời gian thực}
    \label{fig:ad01}
\end{figure}

Quy trình cảnh báo vi phạm vùng an toàn (Hình \ref{fig:ad02}) mô tả luồng xử lý khi thiết bị gửi tọa độ nằm ngoài geofence: backend phát hiện vi phạm, tạo bản ghi cảnh báo, gửi thông báo qua Telegram và đẩy cảnh báo đến ứng dụng di động.

\begin{figure}[H]
    \centering
    \includegraphics[width=0.85\textwidth]{Hinhve/activity_geofence.png}
    \caption{AD02 – Biểu đồ hoạt động: Cảnh báo vi phạm vùng an toàn}
    \label{fig:ad02}
\end{figure}

Quy trình xử lý tín hiệu SOS (Hình \ref{fig:ad03}) thể hiện luồng khẩn cấp khi trẻ nhấn nút SOS trên thiết bị: thiết bị gửi ngay tín hiệu SOS kèm tọa độ, backend tạo cảnh báo khẩn và thông báo đến phụ huynh qua tất cả các kênh đã cấu hình.

\begin{figure}[H]
    \centering
    \includegraphics[width=0.85\textwidth]{Hinhve/activity_sos.png}
    \caption{AD03 – Biểu đồ hoạt động: Gửi tín hiệu SOS khẩn cấp}
    \label{fig:ad03}
\end{figure}

\section{Đặc tả chức năng}
\label{section:2.3}

\subsection{Đặc tả use case: Đăng ký tài khoản}

\begin{table}[H]
\centering
\caption{Đặc tả use case Đăng ký tài khoản}
\label{table:uc_register}
\begin{tabular}{|l|p{9.5cm}|}
\hline
\textbf{Tên use case}       & Đăng ký tài khoản \\
\hline
\textbf{Tác nhân}           & Phụ huynh \\
\hline
\textbf{Tiền điều kiện}     & Phụ huynh chưa có tài khoản trên hệ thống. \\
\hline
\textbf{Luồng sự kiện chính} &
1. Phụ huynh chọn chức năng ``Đăng ký'' trên màn hình chào. \newline
2. Phụ huynh nhập địa chỉ email, mật khẩu và xác nhận mật khẩu. \newline
3. Hệ thống kiểm tra email chưa được đăng ký và mật khẩu đủ độ phức tạp. \newline
4. Hệ thống lưu tài khoản với mật khẩu được băm bằng bcrypt. \newline
5. Ứng dụng chuyển tự động sang màn hình đăng nhập. \\
\hline
\textbf{Luồng phát sinh}    & Nếu địa chỉ email đã tồn tại trong hệ thống, hệ thống thông báo lỗi và yêu cầu nhập địa chỉ khác. Nếu hai trường mật khẩu không khớp nhau, ứng dụng hiển thị thông báo ngay trên giao diện mà không gửi yêu cầu lên server. \\
\hline
\textbf{Hậu điều kiện}      & Tài khoản mới được tạo thành công; phụ huynh có thể đăng nhập và bắt đầu sử dụng hệ thống. \\
\hline
\end{tabular}
\end{table}

\subsection{Đặc tả use case: Đăng ký thiết bị GPS}

\begin{table}[H]
\centering
\caption{Đặc tả use case Đăng ký thiết bị GPS}
\label{table:uc_device}
\begin{tabular}{|l|p{9.5cm}|}
\hline
\textbf{Tên use case}       & Đăng ký thiết bị GPS \\
\hline
\textbf{Tác nhân}           & Phụ huynh \\
\hline
\textbf{Tiền điều kiện}     & Phụ huynh đã đăng nhập; thiết bị IoT đã được lắp đặt và có mã thiết bị (device ID). \\
\hline
\textbf{Luồng sự kiện chính} &
1. Phụ huynh chọn chức năng ``Thêm thiết bị'' trong ứng dụng. \newline
2. Phụ huynh nhập mã thiết bị (device ID) được ghi trên thiết bị IoT. \newline
3. Hệ thống kiểm tra mã thiết bị chưa được đăng ký bởi tài khoản khác. \newline
4. Hệ thống liên kết thiết bị với tài khoản phụ huynh với quyền sở hữu đầy đủ. \newline
5. Ứng dụng hiển thị thiết bị mới trong danh sách và bắt đầu nhận dữ liệu vị trí từ thiết bị đó. \\
\hline
\textbf{Luồng phát sinh}    & Nếu mã thiết bị đã được đăng ký bởi tài khoản khác, hệ thống thông báo ``Thiết bị đã thuộc về tài khoản khác''. Phụ huynh có thể chia sẻ quyền theo dõi với thành viên khác trong gia đình bằng cách nhập email của họ; người được chia sẻ có thể xem vị trí nhưng không thể xóa hay chia sẻ tiếp. \\
\hline
\textbf{Hậu điều kiện}      & Thiết bị được liên kết với tài khoản phụ huynh; dữ liệu vị trí từ thiết bị bắt đầu được lưu và hiển thị trên bản đồ. \\
\hline
\end{tabular}
\end{table}

\subsection{Đặc tả use case: Đăng nhập hệ thống}

\begin{table}[H]
\centering
\caption{Đặc tả use case Đăng nhập hệ thống}
\label{table:uc_login}
\begin{tabular}{|l|p{9.5cm}|}
\hline
\textbf{Tên use case}       & Đăng nhập hệ thống \\
\hline
\textbf{Tác nhân}           & Phụ huynh \\
\hline
\textbf{Tiền điều kiện}     & Phụ huynh đã có tài khoản đăng ký trên hệ thống. \\
\hline
\textbf{Luồng sự kiện chính} &
1. Phụ huynh nhập địa chỉ email và mật khẩu vào ứng dụng. \newline
2. Hệ thống xác thực thông tin đăng nhập; nếu hợp lệ trả về token phiên làm việc. \newline
3. Ứng dụng lưu token và chuyển đến màn hình chính. \\
\hline
\textbf{Luồng phát sinh}    & Nếu thông tin không hợp lệ, hệ thống hiển thị thông báo lỗi. Nếu phiên làm việc hết hạn, ứng dụng tự động chuyển về màn hình đăng nhập. \\
\hline
\textbf{Hậu điều kiện}      & Phụ huynh đã đăng nhập thành công và có thể truy cập toàn bộ chức năng của hệ thống. \\
\hline
\end{tabular}
\end{table}

\subsection{Đặc tả use case: Xem vị trí trẻ theo thời gian thực}

\begin{table}[H]
\centering
\caption{Đặc tả use case Xem vị trí theo thời gian thực}
\label{table:uc_realtime}
\begin{tabular}{|l|p{9.5cm}|}
\hline
\textbf{Tên use case}       & Xem vị trí trẻ theo thời gian thực \\
\hline
\textbf{Tác nhân}           & Phụ huynh, Thiết bị IoT \\
\hline
\textbf{Tiền điều kiện}     & Phụ huynh đã đăng nhập; thiết bị GPS đang hoạt động và có tín hiệu GPS hợp lệ. \\
\hline
\textbf{Luồng sự kiện chính} &
1. Phụ huynh mở màn hình bản đồ trong ứng dụng. \newline
2. Ứng dụng thiết lập kết nối WebSocket đến backend và đăng ký kênh theo dõi thiết bị. \newline
3. Thiết bị IoT gửi tọa độ GPS lên topic MQTT mỗi 5 giây. \newline
4. Backend nhận tọa độ, lưu vào cơ sở dữ liệu và chuyển tiếp đến ứng dụng qua WebSocket. \newline
5. Ứng dụng cập nhật vị trí marker trên bản đồ theo mỗi bản tin nhận được. \\
\hline
\textbf{Luồng phát sinh}    & Nếu thiết bị mất tín hiệu GPS (ở trong nhà, hầm ngầm), marker hiển thị vị trí cuối cùng đã biết kèm nhãn ``Mất tín hiệu'' và thời điểm cập nhật cuối. Nếu kết nối WebSocket bị ngắt, ứng dụng tự động kết nối lại sau 3 giây. \\
\hline
\textbf{Hậu điều kiện}      & Bản đồ hiển thị vị trí mới nhất của trẻ và tiếp tục cập nhật theo thời gian thực. \\
\hline
\end{tabular}
\end{table}

\subsection{Đặc tả use case: Tạo vùng an toàn địa lý}

\begin{table}[H]
\centering
\caption{Đặc tả use case Tạo vùng an toàn địa lý}
\label{table:uc_geofence}
\begin{tabular}{|l|p{9.5cm}|}
\hline
\textbf{Tên use case}       & Tạo vùng an toàn địa lý (Geofence) \\
\hline
\textbf{Tác nhân}           & Phụ huynh \\
\hline
\textbf{Tiền điều kiện}     & Phụ huynh đã đăng nhập. \\
\hline
\textbf{Luồng sự kiện chính} &
1. Phụ huynh chọn chức năng tạo vùng an toàn trong ứng dụng. \newline
2. Phụ huynh chọn loại vùng an toàn: \newline
\hspace{1em}2a. \textit{Vùng tròn}: Phụ huynh nhấn chọn điểm trung tâm trên bản đồ và kéo thanh trượt để điều chỉnh bán kính; hệ thống cảnh báo khi trẻ vượt ra ngoài vùng tròn này. \newline
\hspace{1em}2b. \textit{Đường đi}: Phụ huynh chọn lần lượt các điểm trên bản đồ để vẽ hành trình dự kiến; hệ thống cảnh báo khi trẻ lệch khỏi hành lang cho phép xung quanh đường đi. \newline
3. Phụ huynh đặt tên cho vùng an toàn và xác nhận. \newline
4. Backend lưu cấu hình và phát sự kiện qua WebSocket; vùng an toàn hiển thị ngay trên bản đồ. \\
\hline
\textbf{Luồng phát sinh}    & Nếu phụ huynh chọn chế độ đường đi nhưng chỉ chọn được một điểm (cần tối thiểu 2 điểm để tạo đoạn đường), hệ thống không ghi nhận vùng an toàn hợp lệ; phụ huynh cần chọn thêm điểm hoặc chuyển sang chế độ vùng tròn. \\
\hline
\textbf{Hậu điều kiện}      & Geofence được lưu thành công; backend bắt đầu kiểm tra vi phạm theo cấu hình mới với mỗi điểm vị trí nhận được. \\
\hline
\end{tabular}
\end{table}

\subsection{Đặc tả use case: Cảnh báo vi phạm vùng an toàn}

\begin{table}[H]
\centering
\caption{Đặc tả use case Cảnh báo vi phạm vùng an toàn}
\label{table:uc_alert}
\begin{tabular}{|l|p{9.5cm}|}
\hline
\textbf{Tên use case}       & Cảnh báo vi phạm vùng an toàn \\
\hline
\textbf{Tác nhân}           & Thiết bị IoT, Dịch vụ thông báo, Phụ huynh \\
\hline
\textbf{Tiền điều kiện}     & Có ít nhất một geofence đang hoạt động; thiết bị IoT đang gửi dữ liệu định kỳ. \\
\hline
\textbf{Luồng sự kiện chính} &
1. Hệ thống nhận vị trí mới từ thiết bị. \newline
2. Hệ thống kiểm tra vị trí có nằm trong vùng an toàn đang hoạt động không. \newline
3. Nếu trẻ đã ra ngoài vùng an toàn, hệ thống tạo bản ghi cảnh báo và gửi thông báo qua Telegram đến phụ huynh. \newline
4. Hệ thống đẩy thông báo cảnh báo đến ứng dụng của phụ huynh. \newline
5. Ứng dụng hiển thị cảnh báo với thông tin vị trí hiện tại và tên vùng an toàn bị vi phạm. \\
\hline
\textbf{Luồng phát sinh}    & Nếu trẻ nhanh chóng quay lại vùng an toàn, hệ thống gửi thêm thông báo ``Đã trở về vùng an toàn''. Để tránh spam, hệ thống áp dụng thời gian chờ tối thiểu 60 giây giữa hai cảnh báo liên tiếp cho cùng một geofence. \\
\hline
\textbf{Hậu điều kiện}      & Phụ huynh nhận được cảnh báo qua ít nhất một kênh thông báo. \\
\hline
\end{tabular}
\end{table}

\subsection{Đặc tả use case: Gửi tín hiệu SOS khẩn cấp}

\begin{table}[H]
\centering
\caption{Đặc tả use case Gửi tín hiệu SOS khẩn cấp}
\label{table:uc_sos}
\begin{tabular}{|l|p{9.5cm}|}
\hline
\textbf{Tên use case}       & Gửi tín hiệu SOS khẩn cấp \\
\hline
\textbf{Tác nhân}           & Thiết bị IoT, Dịch vụ thông báo, Phụ huynh \\
\hline
\textbf{Tiền điều kiện}     & Thiết bị IoT đang hoạt động và có kết nối MQTT đến server. \\
\hline
\textbf{Luồng sự kiện chính} &
1. Trẻ em nhấn nút SOS trên thiết bị. \newline
2. Thiết bị đọc tọa độ GPS hiện tại và gửi tín hiệu SOS kèm vị trí lên hệ thống. \newline
3. Hệ thống nhận tín hiệu SOS và tạo bản ghi cảnh báo khẩn. \newline
4. Hệ thống gửi ngay thông báo SOS qua Telegram đến phụ huynh kèm vị trí và đường link bản đồ. \newline
5. Hệ thống đẩy cảnh báo đến ứng dụng của phụ huynh. \newline
6. Ứng dụng hiển thị popup cảnh báo khẩn và vị trí hiện tại trên bản đồ. \\
\hline
\textbf{Luồng phát sinh}    & Nếu thiết bị chưa có tín hiệu GPS tại thời điểm nhấn nút, bản tin SOS vẫn được gửi lên kèm tọa độ cuối cùng đã biết và ghi chú ``Vị trí có thể không chính xác''. Nếu kết nối MQTT bị mất, thiết bị lưu sự kiện SOS vào bộ nhớ và gửi lại ngay khi kết nối được khôi phục. \\
\hline
\textbf{Hậu điều kiện}      & Phụ huynh nhận được cảnh báo SOS qua tất cả các kênh thông báo đã cấu hình; sự kiện SOS được lưu vào lịch sử để tra cứu sau. \\
\hline
\end{tabular}
\end{table}

\subsection{Đặc tả use case: Xem thống kê hành trình}

\begin{table}[H]
\centering
\caption{Đặc tả use case Xem thống kê hành trình}
\label{table:uc_stats}
\begin{tabular}{|l|p{9.5cm}|}
\hline
\textbf{Tên use case}       & Xem thống kê hành trình \\
\hline
\textbf{Tác nhân}           & Phụ huynh \\
\hline
\textbf{Tiền điều kiện}     & Phụ huynh đã đăng nhập; hệ thống đã có dữ liệu vị trí và cảnh báo được lưu trữ. \\
\hline
\textbf{Luồng sự kiện chính} &
1. Phụ huynh chọn chức năng thống kê và chọn khoảng thời gian cần xem. \newline
2. Hệ thống truy vấn dữ liệu vị trí và cảnh báo trong khoảng thời gian đã chọn. \newline
3. Ứng dụng hiển thị các chỉ số tổng hợp: tổng quãng đường di chuyển và số lần ra khỏi vùng an toàn. \newline
4. Ứng dụng vẽ hành trình trên bản đồ thể hiện toàn bộ đường đi của trẻ trong khoảng thời gian đó. \\
\hline
\textbf{Luồng phát sinh}    & Nếu không có dữ liệu cho khoảng thời gian đã chọn, ứng dụng hiển thị các chỉ số bằng 0 và bản đồ trống. \\
\hline
\textbf{Hậu điều kiện}      & Phụ huynh nắm được tổng quãng đường trẻ đã di chuyển, số lần vi phạm vùng an toàn, và hành trình cụ thể trên bản đồ. \\
\hline
\end{tabular}
\end{table}

\section{Yêu cầu phi chức năng}
\label{section:2.4}

\textbf{Hiệu năng.} Hệ thống phải đảm bảo khả năng cập nhật vị trí của thiết bị lên ứng dụng trong thời gian ngắn nhằm đáp ứng yêu cầu theo dõi theo thời gian thực. Backend cần duy trì khả năng tiếp nhận, xử lý và lưu trữ dữ liệu vị trí ổn định khi có nhiều thiết bị kết nối đồng thời.

\textbf{Độ tin cậy.} Thiết bị phải có khả năng tự động khôi phục kết nối với máy chủ khi xảy ra gián đoạn mạng mà không cần sự can thiệp của người dùng. Hệ thống phải duy trì khả năng theo dõi liên tục và hạn chế tối đa việc mất dữ liệu trong quá trình truyền nhận.

\textbf{Tính dễ sử dụng.} Giao diện ứng dụng di động cần được thiết kế trực quan, dễ sử dụng và phù hợp với người dùng không có nhiều kinh nghiệm về công nghệ. Các chức năng chính phải được bố trí hợp lý để người dùng có thể thao tác nhanh chóng.

\textbf{Bảo mật.} Hệ thống phải xác thực người dùng trước khi cho phép truy cập các chức năng và dữ liệu. Thông tin đăng nhập cần được lưu trữ an toàn nhằm ngăn chặn nguy cơ rò rỉ hoặc truy cập trái phép. Dữ liệu vị trí GPS trong quá trình truyền giữa thiết bị IoT, backend và ứng dụng di động phải được bảo vệ bằng cơ chế mã hóa để đảm bảo tính bảo mật, tính toàn vẹn dữ liệu và hạn chế nguy cơ bị nghe lén hoặc giả mạo trong quá trình truyền nhận. Đồng thời, chỉ những người dùng được cấp quyền mới được phép truy cập và quản lý dữ liệu của thiết bị.

\textbf{Yêu cầu kỹ thuật.} Hệ thống phải được xây dựng theo mô hình client--server, trong đó backend được triển khai trên nền tảng điện toán đám mây nhằm đảm bảo khả năng truy cập từ Internet cho cả thiết bị IoT và ứng dụng di động. Kiến trúc hệ thống cần hỗ trợ khả năng mở rộng khi số lượng thiết bị tăng lên, đồng thời hỗ trợ lưu trữ dữ liệu vị trí để phục vụ việc theo dõi và tra cứu lịch sử di chuyển.

Kết chương: Chương 2 đã phân tích chi tiết hiện trạng các giải pháp theo dõi vị trí trẻ em, chỉ ra các hạn chế còn tồn tại và xác định rõ yêu cầu cho hệ thống cần xây dựng. Kết quả khảo sát cho thấy hướng tiếp cận thiết bị IoT chuyên dụng kết hợp GPS và kết nối di động là phù hợp nhất. Các yêu cầu chức năng và phi chức năng được đặc tả trong chương này sẽ là cơ sở để lựa chọn công nghệ và thiết kế kiến trúc hệ thống, được trình bày trong các chương tiếp theo.

\end{document}
